\documentclass[UTF8,zihao=-4]{ctexart}
\usepackage[a4paper,margin=2.1cm]{geometry}
\usepackage{hyperref}
\usepackage{array}
\usepackage{longtable}
\usepackage{verbatim}
\hypersetup{hidelinks}
\setlength{\parindent}{0pt}
\setlength{\parskip}{0.55em}
\emergencystretch=4em
\sloppy
\title{课程框架与整理方法}
\author{}
\date{}
\begin{document}
\maketitle
\setcounter{tocdepth}{1}
\tableofcontents
\newpage
\section{编译原理零基础突击路线图}


\subsection{目标}


这套讲义按“期末能写出步骤分”整理。每一章只保留四类内容：


\begin{itemize}
\item 必背定义：题目要求“解释、说明、定义”时直接写。
\item 固定算法：自动机、LL、LR、回填、SSA、数据流、Andersen。
\item 最小例子：看懂题型，不展开长推导。
\item 真题答法：按题目要求列步骤、写表、写结论。
\end{itemize}


阅读顺序：


\begin{quote}
\footnotesize
\begin{verbatim}
01 词法分析 -> 02 语法分析 -> 03 语义与中间代码
-> 04 后端 -> 05 程序分析 -> Exam/答案整理
\end{verbatim}
\end{quote}


\subsection{分数优先级}


第一优先级，先会做题：


\begin{itemize}
\item 正则表达式、Thompson NFA、子集构造 DFA、DFA 最小化。
\item 消除左递归、FIRST、FOLLOW、预测分析表、LL(1) 判断。
\item 语法制导定义、三地址码、布尔表达式回填。
\item 指令调度三类依赖、寄存器重命名、列表调度。
\item 支配、支配边界、SSA 中 phi 的放置位置。
\item 符号执行、CNF、Tseitin、Andersen 指针分析。
\end{itemize}


第二优先级，掌握模板：


\begin{itemize}
\item LR(0)、SLR(1)、LR(1) 的项目、closure、goto、分析表。
\item ANTLR4 左递归表达式改写与优先级上升。
\item 寄存器分配、SSA 转出、基本数据流分析。
\end{itemize}


第三优先级，背定义拿小分：


\begin{itemize}
\item 编译器阶段划分。
\item CFG、NPDA、CFL 封闭性。
\item flow-sensitive、path-sensitive、context-sensitive、field-sensitive。
\end{itemize}


\subsection{两天突击安排}


第一天：


\begin{quote}
\footnotesize
\begin{verbatim}
上午：01 词法分析，重点手算 NFA -> DFA -> 最小化
下午：02 语法分析，重点 LL(1)，再看 LR 模板
晚上：03 语义与中间代码，重点 SDD 与回填
\end{verbatim}
\end{quote}


第二天：


\begin{quote}
\footnotesize
\begin{verbatim}
上午：04 后端，重点调度、依赖、寄存器分配定义
下午：05 程序分析，重点数据流、符号执行、Andersen
晚上：Exam/答案整理，按年份过答题模板
\end{verbatim}
\end{quote}


\subsection{答题格式}


自动机题：


\begin{quote}
\footnotesize
\begin{verbatim}
1. 按正则结构分层。
2. 套 Thompson 模板画 NFA。
3. 写 epsilon-closure 和 move。
4. 列 DFA 状态集合表。
5. 标终态。
6. 按终态/非终态初分组，迭代拆分完成最小化。
\end{verbatim}
\end{quote}


LL 题：


\begin{quote}
\footnotesize
\begin{verbatim}
1. 消除左递归。
2. 写 FIRST。
3. 写 FOLLOW。
4. 填预测分析表。
5. 表中无多重入口就是 LL(1)，有多重入口就不是。
\end{verbatim}
\end{quote}


中间代码题：


\begin{quote}
\footnotesize
\begin{verbatim}
1. 写 truelist、falselist、nextlist。
2. 为条件跳转留空目标。
3. 用 backpatch 填循环入口、循环体入口、循环出口。
4. 最后写出带编号的三地址码。
\end{verbatim}
\end{quote}


分析题：


\begin{quote}
\footnotesize
\begin{verbatim}
先写定义，再写构造规则，最后写结论。
\end{verbatim}
\end{quote}


\subsection{文件说明}


\begin{itemize}
\item \texttt{01-编译器概览与词法分析复习讲义.md/.pdf}：正则、NFA、DFA、最小化。
\item \texttt{02-语法分析复习讲义.md/.pdf}：CFG、LL、LR、ANTLR。
\item \texttt{03-语义分析与中间代码复习讲义.md/.pdf}：SDD、符号表、TAC、回填、SSA。
\item \texttt{04-代码生成与优化复习讲义.md/.pdf}：目标代码、寄存器分配、指令调度。
\item \texttt{05-程序分析与符号执行复习讲义.md/.pdf}：数据流、符号执行、SAT、Andersen。
\item \texttt{Exam/答案整理/期末试题答案汇总.md/.pdf}：真题题型和答题模板。
\end{itemize}

\end{document}
